\subsubsection{Scaling the Marker Results to Identify Troubled Cells}
\label{SEC:ScaledMarker}
\key*{}[marker!scaled marker]

%%%%%%%%%%%%%%%%%%%%%%%%%%%%%%%%%%%%%%%%%%%%%%%%%%%%%%%%%%%%%%%%%%%%%%%%%%%%%%%%

From the local marker, a local {\em scaled} marker can be computed
\begin{align}
    \label{EQN:ScaledResolutionMarker}
    \scaledmarker
 =& \begin{cases}
      0                                  &
      \phantom{\bigl|} \zeta \phantom{\bigr|} \leq -\pi/2
 \\
      \left( 1 + \sin \zeta \right) / 2  & 
               \bigl|  \zeta          \bigr|  \leq  \pi/2
 \\
      1                                  & 
      \phantom{\bigl|} \zeta \phantom{\bigr|} \geq  \pi/2
    \end{cases}
  & 
    \zeta
 =& \frac{\pi}{2} \, \frac{ \marker - \Psi_0 }{ \Delta \Psi }
\end{align}
where
\begin{align}
    \Psi_0 =& -\left[ C_1 + C_2 \log_{10} (\PO) \right]
  & \Delta \Psi =& 0.5
\end{align}

%%%%%%%%%%%%%%%%%%%%%%%%%%%%%%%%%%%%%%%%%%%%%%%%%%%%%%%%%%%%%%%%%%%%%%%%%%%%%%%%

As shown in \fref{FIG:Scaled_Marker}, the scaled marker varies between
0 (region to retain) and 1 (region to filter).
%
Other works sometimes use this smooth variation to modify the
strength of the dissipation applied.
%
In GFR, this smoothly varying scaled marker can be output for
visualization purposes.
%
However, troubled cells are simply identified by
${\cal S} > \scinum{1.0}{-8}$.
%
This approach is more aggressive at identifying troubled cells.

%%%%%%%%%%%%%%%%%%%%%%%%%%%%%%%%%%%%%%%%%%%%%%%%%%%%%%%%%%%%%%%%%%%%%%%%%%%%%%%%

\begin{figure}
  \begin{centering}
    \insertimage[width=4.0in]{Scaled_Marker/Scaled_Marker}
    \caption{Scaled Marker}
    \label{FIG:Scaled_Marker}
  \end{centering}
\end{figure}

%%%%%%%%%%%%%%%%%%%%%%%%%%%%%%%%%%%%%%%%%%%%%%%%%%%%%%%%%%%%%%%%%%%%%%%%%%%%%%%%

%% Description for the old version, when $(C_1,C_2)$ were still inputs:
%
% The default values of the $(C_1,C_2)$ coefficients are
% (4.00,4.25) when using the resolution marker and
% (2.25,3.00) when using the jump marker.
% %
% GFR allows one to either specify $(C_1,C_2,\Delta\Psi)$
% or to directly specify
% \inp[nmlLimiter::MarkerMin]{Marker\_Min} = $\Psi_{min}$
%                                          = $(\Psi_0-\Delta\Psi)$ and
% \inp[nmlLimiter::MarkerMax]{Marker\_Max} = $\Psi_{max}$
%                                          = $(\Psi_0+\Delta\Psi)$
% in which the variation in the scaled marker $\scaledmarker$ occurs.

The default values of the $(C_1,C_2)$ coefficients are
$(4.0,4.25)$ when using the resolution marker and
$(2.25,3.0)$ when using the jump marker.
%
GFR will use $(C_1,C_2,\Delta\Psi)$ to compute 
$\Psi_{min} = \Psi_0-\Delta\Psi$ and
$\Psi_{max} = \Psi_0+\Delta\Psi$
unless the inputs for
(\inp[nmlLimiter::MarkerMin]{Marker\_Min},%
 \inp[nmlLimiter::MarkerMax]{Marker\_Max})
=
($\Psi_{min},\Psi_{max}$) are used instead.

%%%%%%%%%%%%%%%%%%%%%%%%%%%%%%%%%%%%%%%%%%%%%%%%%%%%%%%%%%%%%%%%%%%%%%%%%%%%%%%%

When using the MLP marker, the scaled marker is simply computed as
a binary flag to indicate a troubled cell, \thatis any cell in which
the multidimensional limiting process would reduce the solution
polynomial order.
\begin{align}
    \scaledmarker
 =& \begin{cases}
       0 & \text{if $\marker = \PO$}
    \\ 1 & \text{otherwise}
    \end{cases}
\end{align}

%%%%%%%%%%%%%%%%%%%%%%%%%%%%%%%%%%%%%%%%%%%%%%%%%%%%%%%%%%%%%%%%%%%%%%%%%%%%%%%%

\SeeAlso{
\inp[nmlLimiter::MarkerDpsi]{Marker\_Dpsi},
\inp[nmlLimiter::MarkerMin]{Marker\_Min},
\inp[nmlLimiter::MarkerMax]{Marker\_Max}
}

